\chapter{Conclusão}
\label{cap:conclusao}

Este trabalho partiu da constatação de que modelos preditivos treinados com microdados do ENEM, embora aparentemente neutros e de desempenho aceitável, podem converter desigualdades socioeconômicas e raciais em critérios automatizados de exclusão, sem que as métricas convencionais de avaliação revelem o problema. Para enfrentá-lo, foi projetada, desenvolvida e avaliada, segundo a \textit{Design Science Research}, a ferramenta AuditEdu-XAI. A ferramenta integra métricas de justiça algorítmica e explicabilidade baseada em valores de Shapley em uma arquitetura desacoplada, executável em infraestrutura gratuita.

% ---------------------------------------------------------------------------
\section{Atendimento aos Objetivos}

Quanto à questão de pesquisa formulada na Introdução, a resposta é afirmativa, com as ressalvas apresentadas na Seção~\ref{sec:res_limitacoes}. A ferramenta, executável em infraestrutura gratuita, detectou, quantificou e explicou vieses de raça/cor e renda em modelos preditivos do ENEM e mostrou-se sensível a vieses aos quais a AUC-ROC pouco respondeu. O objetivo geral foi atingido: a ferramenta detecta, quantifica e explica vieses em modelos preditivos de desempenho educacional. Os objetivos específicos foram cumpridos da seguinte forma:
\begin{itemize}
    \item \textbf{Pipeline de dados:} os módulos de ingestão e pré-processamento processam os microdados de três edições em blocos e persistem os resultados em Parquet, com leitura em blocos projetada para o limite de 8~GB (RNF02) (Capítulo~\ref{cap:desenvolvimento});
    \item \textbf{Modelo-alvo:} classificadores XGBoost foram otimizados por validação cruzada e avaliados em conjunto de teste isolado, com AUC-ROC em torno de 0,75 (Seção~\ref{sec:res_desempenho});
    \item \textbf{Motor de equidade:} o artefato calcula Paridade Demográfica, Impacto Disparate, DI dos rótulos reais, TPR e FPR exclusivamente sobre dados não vistos no treinamento (Seções~\ref{sec:res_raca} a~\ref{sec:res_amplificacao});
    \item \textbf{Explicabilidade:} a camada TreeSHAP quantifica a contribuição de cada atributo, global e localmente (Seção~\ref{sec:res_shap});
    \item \textbf{Interface:} o painel Streamlit atende aos requisitos funcionais, conforme o Quadro~\ref{qua:req_evidencia}, com tempo de resposta inferior a 2~s em todas as execuções medidas em ambiente local, inclusive com o processamento restrito a 2 núcleos (Seção~\ref{sec:desempenho});
    \item \textbf{Validação:} o experimento de injeção sintética de viés corroborou as três hipóteses: sensibilidade, especificidade e baixa sensibilidade das métricas convencionais (Seção~\ref{sec:res_injecao}).
\end{itemize}

% ---------------------------------------------------------------------------
\section{Contribuições}

As contribuições do trabalho situam-se em três dimensões.

\textbf{Contribuição tecnológica.} Um artefato, com código publicado em repositório público, que combina, em um mesmo fluxo, métricas de equidade de grupo, decomposição do viés entre dados e modelo, explicações TreeSHAP e uma interface acessível a usuários sem formação em programação. Sua arquitetura desacoplada, com pré-cálculo em lote e leitura de arquivos estáticos, demonstra, com a implantação pública em plataforma gratuita (\url{https://auditeduxai.streamlit.app/}), que a auditoria algorítmica pode ser oferecida sem custo de infraestrutura. Além disso, a imagem Docker validada por integração contínua assegura a portabilidade para outros ambientes. O padrão é replicável para outras bases públicas.

\textbf{Contribuição empírica.} A aplicação aos dados de 2019 a 2021 produziu três achados substantivos:
\begin{itemize}
    \item um modelo com desempenho moderado viola a regra dos quatro quintos para todos os grupos não brancos e para candidatos com renda familiar de até 2,5 salários mínimos, nas três edições;
    \item o modelo \textbf{amplifica} a disparidade presente nos dados, de forma estatisticamente robusta em 13 das 15 combinações de raça/cor e edição e nas faixas de renda mais baixas, levando para baixo do limiar, de forma robusta, grupos cujo DI nos rótulos reais estava na fronteira de 0,8;
    \item a ausência de computador na residência é o atributo mais influente nas predições, e sua importância cresce entre 2019 e 2021, em coincidência temporal com o ensino remoto da pandemia. O achado é compatível com a hipótese de agravamento da exclusão digital, embora não permita inferência causal.
\end{itemize}

\textbf{Contribuição metodológica.} Um protocolo de validação por injeção de viés com controle negativo pareado e efeito líquido, que separa o efeito do viés direcionado do efeito genérico de alterar a prevalência de rótulos. O protocolo evidenciou que, sob um viés que leva o DI de um grupo a praticamente zero, a AUC-ROC varia, em média, no máximo 0,032 e, para o grupo Indígena, menos de 0,001. Esse resultado mostra concretamente por que a avaliação de modelos aplicados a populações heterogêneas não pode se limitar a métricas agregadas de desempenho.

% ---------------------------------------------------------------------------
\section{Limitações}

As limitações, detalhadas na Seção~\ref{sec:res_limitacoes}, podem ser resumidas em nove pontos:
\begin{itemize}
    \item o pequeno número de candidatos indígenas no conjunto de teste, que torna suas métricas instáveis;
    \item a auditoria de um único tipo de modelo, que usa o atributo raça/cor diretamente;
    \item a natureza associativa, e não causal, dos valores SHAP;
    \item a dicotomização do desempenho pela mediana;
    \item a comparabilidade limitada entre edições, decorrente da variação do salário mínimo e da composição dos participantes;
    \item o alvo relativo a cada edição, que expressa a posição do candidato em relação à mediana do ano, e não um patamar absoluto de proficiência;
    \item a imputação pela moda, que em 2020 atribui um perfil presumido a participantes cuja ausência de resposta provavelmente não é aleatória;
    \item a seleção de hiperparâmetros na borda de uma grade compacta;
    \item a ausência de medição do tempo de resposta e do consumo de memória no ambiente de produção.
\end{itemize}

% ---------------------------------------------------------------------------
\section{Trabalhos Futuros}

Os resultados abrem as seguintes linhas de continuidade:
\begin{itemize}
    \item \textbf{Incerteza estatística:} estender os intervalos de confiança, que hoje refletem apenas a variabilidade do conjunto de teste, a um \textit{bootstrap} com retreinamento \cite{efron1993bootstrap}, e ampliar a amostra dos grupos minoritários por sobreamostragem estratificada;
    \item \textbf{Tratamento de não respostas:} substituir a imputação pela moda por uma categoria explícita de ``não resposta'' ao questionário;
    \item \textbf{Validação em produção:} medir o tempo de resposta e o consumo de memória no próprio ambiente de implantação;
    \item \textbf{Viés por variáveis \textit{proxy}:} auditar um modelo treinado \textbf{sem} o atributo raça/cor, para medir quanto da disparidade persiste por meio de variáveis correlacionadas \cite{dwork2012fairness};
    \item \textbf{Mitigação:} integrar técnicas de pré-processamento, como a reponderação de instâncias \cite{kamiran2012data}, e de pós-processamento, como limiares por grupo \cite{hardt2016equality}, e exibir o \textit{trade-off} entre desempenho e equidade na própria interface;
    \item \textbf{Interseccionalidade:} avaliar subgrupos definidos pela combinação de atributos, como raça/cor e renda, nos quais disparidades podem ser maiores que as observadas em cada atributo isolado \cite{buolamwini2018gender};
    \item \textbf{Generalização:} estender a auditoria a outras famílias de modelos, a grades de hiperparâmetros mais amplas, a outros critérios de definição do alvo, a novas edições do ENEM e a outras bases educacionais públicas, como o SAEB e o Censo Escolar;
    \item \textbf{Interface:} avaliar a usabilidade do painel com usuários do público-alvo, como gestores e membros de órgãos de controle, por meio de testes de usabilidade.
\end{itemize}

% ---------------------------------------------------------------------------
\section{Considerações Finais}

A adoção de modelos preditivos na gestão educacional é uma tendência, e a regulação brasileira caminha para exigir dos sistemas de alto risco avaliações de impacto algorítmico e mecanismos de explicação \cite{brasil2018lgpd, brasil2023pl2338}. Os resultados deste trabalho reforçam a pertinência de tais exigências. Um modelo pode parecer adequado por todas as métricas usuais e, ainda assim, excluir sistematicamente os grupos que mais dependem de políticas educacionais justas. A ferramenta desenvolvida torna essa exclusão visível, mensurável e explicável, e, ao fazê-lo, oferece a gestores, pesquisadores e órgãos de controle um instrumento concreto para que a promessa de decisões orientadas a dados não se converta em desigualdade automatizada.
